\section{SD-C11: the holomorphic reflected double}
\label{sec:source}

Let $\At=\{F_{p_n}\}$ be the entropy-ordered tensor atoms.  Retain the
recurrent nearest-neighbor base grammar of SD-C10.  Let $\mathbb F_+$ and
$\mathbb F_-$ be free groups generated by disjoint directed-positive edge
alphabets.  Their group von Neumann algebras carry canonical traces
$\tau_+$ and $\tau_-$.  Work in
\[
 \cN=B(\ell^2(\At))\bar\otimes
 L(\mathbb F_+)\bar\otimes L(\mathbb F_-),
 \quad \Phi=\Tr\otimes\tau_+\otimes\tau_-.   \tag{3.1}
\]

For $\varepsilon\in\{+,-\}$, let $T_s^\varepsilon$ be the SD-C10 transfer
with its own positive alphabet:
\begin{align}
 T_s^\varepsilon={}&\sum_np_n^{-s}E_{nn}\otimes1\notag\\
 &+\sum_n\frac{p_n^{-s}+p_{n+1}^{-s}}2
 \left(E_{n+1,n}\otimes\lambda(g_{n,+}^\varepsilon)
 +E_{n,n+1}\otimes\lambda(g_{n,-}^\varepsilon)\right). \tag{3.2}
\end{align}
The unused group tensor factor acts as the identity.  Define
\[
 \boxed{
 \cC_s=\begin{pmatrix}0&T_s^+\\T_{1-s}^-&0\end{pmatrix}}
 \quad\text{on }\CC^2\otimes\ell^2(\At)\otimes
 \ell^2(\mathbb F_+\times\mathbb F_-).       \tag{3.3}
\]

The lower block is the same positive-monoid representation evaluated at
$1-s$.  It is not the adjoint, not the group-ring transpose, and does not
apply $g\mapsto g^{-1}$.  This convention is essential: each of those
involutions would restore identity words that the positive cone is meant to
filter.

This freezes SD-C11.  The two channel states form a period-two symbolic
extension; every step switches layers.  Atom-loop weights are $p^{-s}$ in
the plus layer and $p^{-(1-s)}$ in the minus layer.  Cross-edge endpoint roofs
and potentials are inherited unchanged.  The determinant convention is the
noncommutative regularized trace-series determinant $\det_q(I-z\cC_s)$ on a
common $L^q$ strip.

No Riemann zeros, target scales, fitted phases, or external system family are
allowed.  Reflection is the fixed layer swap
\[
 J=\begin{pmatrix}0&I\\I&0\end{pmatrix}.     \tag{3.4}
\]

\begin{proposition}[Common ideal strip and reflection]
\label{prop:strip}
For integer $q\ge1$,
\[
 \cC_s\in L^q\quad\Longleftrightarrow\quad
 \frac1q<\RePart s<1-\frac1q.               \tag{3.5}
\]
The first nonempty integer strip is $q=3$.  Moreover, after exchanging the
two independent alphabets by their canonical generator bijection,
\[
 J\cC_sJ=\cC_{1-s}.                          \tag{3.6}
\]
Consequently every normalized trace-series $\det_q$ obeys
$\Delta_q(s,z)=\Delta_q(1-s,z)$.
\end{proposition}

\begin{proof}
The two off-diagonal blocks belong to $L^q$ precisely when
$q\Re s>1$ and $q(1-\Re s)>1$, as in the banded ideal theorem for SD-C10.
Layer exchange swaps the blocks; the alphabet bijection identifies their
otherwise disjoint universal copies.  Trace and regularized determinant are
invariant under this trace-preserving conjugacy.
\end{proof}

The block-exchange identity holds for any family inserted in this two-channel
template.  It is therefore a proves-too-much control, not by itself a Riemann
functional equation.  Arithmetic credit depends entirely on the filtered
closed-word ledger studied next.
